\section{Explanation by Example}
\label{sec:explanation}

This section previews the slicing apparatus formally introduced in later chapters through three worked examples. \Cref{fig:parse-digit-slices} contrasts \synsrc{synthesis} and \anasrc{analysis} slices on a single function call --- the two judgements are defined in \cref{sec:synthesis,sec:analysis}. \Cref{fig:parse-pin-slices} demonstrates \emph{query refinement} (the formal apparatus appears in \cref{sec:decompositions}) and the use of slicing for static type-error debugging, treated systematically in \cref{sec:marking}. \Cref{fig:parse-event-slices} demonstrates \emph{dynamic sum types}: an unannotated injection \texttt{Left(a)} is inferred at \texttt{A~+~$\dyn$} and \texttt{Right(b)} at \texttt{$\dyn$~+~B}, with $\dyn$ filling the unmentioned alternative and so removing the annotation otherwise required at every injection. The case rule then joins the two branches into a precise sum, e.g.\ \texttt{(A~+~$\dyn$)~$\sqcup$~($\dyn$~+~B)~=~A~+~B} --- the union inference TypeScript performs on tagged-literal branches. Lehmann and Tanter's \emph{Sums of Uncertainty}~\cite{SumsOfUncertainty} develops gradual sum types far further; we take only their two injection rules giving \texttt{Left(a)~:~A~+~$\dyn$} and \texttt{Right(b)~:~$\dyn$~+~B}. In \cref{fig:parse-event-slices} the bug, if any, sits in the selected term itself; equally though, the surrounding context --- here a programmer-supplied annotation --- could have been the wrong party, and type slicing offers symmetric machinery (the analysis slice) for interrogating that side instead.

\paragraph{How to read the figures.}
\begin{itemize}\setlength\itemsep{1pt}
  \item \synsrc{Red} and \anasrc{blue} mark code retained by a synthesis or analysis slice respectively.
  \item \asmref{Orange} marks variable bindings whose own slices are referenced from the current slice as \emph{assumption slices}; their bodies are sliced recursively and shown as nested sub-slices.
  \item \strN{Green} text is \emph{purely structural}: code retained for syntactic well-formedness but not sourcing the queried type. An \texttt{if} or \texttt{case} keyword renders green precisely when the joined type can be sourced entirely from one surviving branch, leaving the construct itself irrelevant to the join.\footnote{Treating a position as purely structural rather than slice-contributing is a presentational refinement, not part of the formal slicing relation; it is straightforward to implement in a presentation tool.}
  \item \protect\fold{} replaces sub-terms outside the slice.
  \item \synsrcS{Bolded} red and \anasrcS{blue} positions are the \emph{type sources} of a slice: a minimal sub-set of slice positions from which the queried type can be reconstructed.\footnote{This bolded sub-slice is informal --- not part of the formal slicing relation --- but is easy to compute in practice.}
\end{itemize}

\begin{figure}[!t]
\centering

% --- Local visual conventions (figure-scoped) ---
% Dice symbol drawn in TikZ so it does not depend on a system font with ⚂ (U+2682).
\providecommand{\fold}{%
  \tikz[baseline=0.1em]{%
    \draw[rounded corners=0.3pt, line width=0.28pt, black!75]
      (0,0) rectangle (0.85em,0.85em);%
    \fill[black!85] (0.18em,0.67em) circle (0.07em);%
    \fill[black!85] (0.425em,0.425em) circle (0.07em);%
    \fill[black!85] (0.67em,0.18em) circle (0.07em);%
  }%
}
% Sources of the queried type: synthesis (red) and analysis (blue).
\providecommand{\synsrc}[1]{\textcolor{BrickRed!60!gray}{#1}}
\providecommand{\anasrc}[1]{\textcolor{NavyBlue!60!gray}{#1}}
\providecommand{\synsrcS}[1]{\textcolor{red!85!black}{\textbf{\boldmath\addfontfeature{FakeBold=4.5}#1}}}
\providecommand{\anasrcS}[1]{\textcolor{blue!85!black}{\textbf{\boldmath\addfontfeature{FakeBold=4.5}#1}}}
% Purely structural skeleton: single middle shade of green.
\providecommand{\strN}[1]{\textcolor{OliveGreen!65!black}{#1}}
\providecommand{\strM}[1]{\textcolor{OliveGreen!65!black}{#1}}
\providecommand{\strF}[1]{\textcolor{OliveGreen!65!black}{#1}}
% Selected (focused) sub-term: soft yellow highlight, no harsh outline.
\providecommand{\focusbox}[1]{%
  \tikz[baseline=(N.base)]{%
    \node[draw=black!30, line width=0.3pt, rounded corners=2pt,
          inner sep=2pt, fill=black!10, anchor=base, align=left] (N)
         {\strut\ttfamily\bfseries #1};%
  }%
}
% Gapped sub-terms shown in their original position (in-context view): muted grey.
\providecommand{\offslice}[1]{\textcolor{black!42}{#1}}

% --- TikZ styles for the code segments and the artistic fold arrow ---
\tikzset{
  codebox/.style={
    draw=black!55, line width=0.5pt, dotted,
    rounded corners=2.5pt, inner sep=5pt,
    align=left, anchor=north,
    font=\ttfamily\scriptsize,
  },
  artyarrow/.style={
    -{Stealth[length=2.8mm, width=2.2mm]},
    line width=0.7pt, color=black!55,
    decorate, decoration={
      snake, amplitude=0.6mm, segment length=2.7mm,
      pre length=2mm, post length=2.5mm,
    },
  },
  arrowlbl/.style={
    font=\itshape\scriptsize, color=black!60,
    fill=white, inner sep=1.5pt,
  },
}

\begin{subfigure}[t]{0.48\textwidth}
\centering
{\textsc{\large\bfseries Synthesis slice}}\par
\smallskip
{\scriptsize\itshape\color{black!55}select }{\scriptsize\synsrc{\texttt{parse\_digit}}}\par
{\scriptsize\itshape\color{black!55}query }{\scriptsize$\upsilon\,=\,$}{\scriptsize\synsrc{\texttt{Char \textrightarrow\ Option Digit}}}
\par\smallskip

\begin{tikzpicture}
\node[codebox] (orig) at (0,0) {%
\strN{type} \asmref{Option}\synsrc{(A)} \strN{=} \synsrc{Some(A)} \strM{|} \offslice{None}\\
\strN{type} \asmref{Digit} \strN{=} \synsrc{Zero} \strM{|} \offslice{One | Two | $\cdots$ | Nine}\\
\strN{let} \asmref{parse\_digit} \strN{=} \synsrcS{fun} \offslice{c} \strM{:} \synsrcS{Char} \synsrcS{-\textgreater}\\
\quad\strM{case} \offslice{c} \strM{of}\\
\qquad\strF{|} \offslice{'0'} \strF{=\textgreater} \synsrcS{Some Zero}\\
\qquad\offslice{| '1' =\textgreater\ Some One}\\
\qquad\offslice{$\;\;\vdots$}\\
\qquad\offslice{| '9' =\textgreater\ Some Nine}\\
\qquad\offslice{| \_\ \ =\textgreater\ None}\\
\strN{in} \focusbox{\synsrc{parse\_digit}}\strN{(}\offslice{'5'}\strN{)}%
};

\node[codebox, below=14mm of orig] (fold) {%
\strN{type} \asmref{Option}\synsrc{(A)} \strN{=} \synsrc{Some(A)} \strM{|} \fold\\
\strN{type} \asmref{Digit} \strN{=} \synsrc{Zero} \strM{|} \fold\\
\strN{let} \asmref{parse\_digit} \strN{=} \synsrcS{fun} \fold{} \strM{:} \synsrcS{Char} \synsrcS{-\textgreater}\\
\quad\strM{case} \fold{} \strM{of}\\
\qquad\strF{|} \fold{} \strF{=\textgreater} \synsrcS{Some Zero}\\
\qquad\strF{|} \fold\\
\strN{in} \focusbox{\synsrc{parse\_digit}}\strN{(}\fold\strN{)}%
};

\draw[artyarrow] (orig.south) -- node[arrowlbl] {fold} (fold.north);

% Enforce a fixed bounding box; identical coords in both subfigures guarantee
% the code segments and captions line up horizontally.
\useasboundingbox (-3.8,0.55) rectangle (3.8,-9.6);
\end{tikzpicture}

\caption{\texttt{parse\_digit} synthesises type \synsrc{\texttt{Char → Option Digit}}.}\label{fig:syn-parse-digit}
\end{subfigure}\hfill
\begin{subfigure}[t]{0.48\textwidth}
\centering
{\textsc{\large\bfseries Analysis slice}}\par
\smallskip
{\scriptsize\itshape\color{black!55}select }{\scriptsize\anasrc{\texttt{'5'}}}\par
{\scriptsize\itshape\color{black!55}query }{\scriptsize$\upsilon\,=\,$}{\scriptsize\anasrc{\texttt{Char}}}
\par\smallskip

\begin{tikzpicture}
\node[codebox] (orig) at (0,0) {%
\offslice{type Option(A) = Some(A) | None}\\
\offslice{type Digit = Zero | One | Two | $\cdots$ | Nine}\\
\strN{let} \asmref{parse\_digit} \strN{=} \anasrc{fun} \offslice{c} \strM{:} \anasrcS{Char} \anasrc{-\textgreater}\\
\quad\offslice{case c of}\\
\qquad\offslice{| '0' =\textgreater\ Some Zero}\\
\qquad\offslice{| '1' =\textgreater\ Some One}\\
\qquad\offslice{$\;\;\vdots$}\\
\qquad\offslice{| '9' =\textgreater\ Some Nine}\\
\qquad\offslice{| \_\ \ =\textgreater\ None}\\
\strN{in} \anasrc{parse\_digit}\anasrc{(}\focusbox{'5'}\anasrc{)}%
};

\node[codebox, below=14mm of orig] (fold) {%
\fold\\
\fold\\
\strN{let} \asmref{parse\_digit} \strN{=} \anasrc{fun} \fold{} \strM{:} \anasrcS{Char} \anasrc{-\textgreater} \fold\\
\strN{in} \anasrc{parse\_digit}\anasrc{(}\focusbox{\fold{}}\anasrc{)}%
};

\draw[artyarrow] (orig.south) -- node[arrowlbl] {fold} (fold.north);

% Same bounding box as the synthesis subfigure → horizontally aligned captions.
\useasboundingbox (-3.8,0.55) rectangle (3.8,-9.6);
\end{tikzpicture}

\caption{\texttt{'5'} analyses against type \anasrcS{\texttt{Char}}.}\label{fig:ana-parse-digit}
\end{subfigure}

\caption{Parsing a character into an \texttt{Option Digit}.}
\label{fig:parse-digit-slices}
\end{figure}

\begin{figure}[p]
\centering

% Local visual conventions (figure-scoped). Most are reused from the previous figure.
\providecommand{\fold}{%
  \tikz[baseline=0.1em]{%
    \draw[rounded corners=0.3pt, line width=0.28pt, black!75]
      (0,0) rectangle (0.85em,0.85em);%
    \fill[black!85] (0.18em,0.67em) circle (0.07em);%
    \fill[black!85] (0.425em,0.425em) circle (0.07em);%
    \fill[black!85] (0.67em,0.18em) circle (0.07em);%
  }%
}
\providecommand{\synsrc}[1]{\textcolor{BrickRed!60!gray}{#1}}
\providecommand{\anasrc}[1]{\textcolor{NavyBlue!60!gray}{#1}}
\providecommand{\synsrcS}[1]{\textcolor{red!85!black}{\textbf{\boldmath\addfontfeature{FakeBold=4.5}#1}}}
\providecommand{\anasrcS}[1]{\textcolor{blue!85!black}{\textbf{\boldmath\addfontfeature{FakeBold=4.5}#1}}}
\providecommand{\strN}[1]{\textcolor{OliveGreen!65!black}{#1}}
\providecommand{\strM}[1]{\textcolor{OliveGreen!65!black}{#1}}
\providecommand{\strF}[1]{\textcolor{OliveGreen!65!black}{#1}}
\providecommand{\offslice}[1]{\textcolor{black!42}{#1}}
\providecommand{\focusbox}[1]{%
  \tikz[baseline=(N.base)]{%
    \node[draw=black!30, line width=0.3pt, rounded corners=2pt,
          inner sep=2pt, fill=black!10, anchor=base, align=left] (N)
         {\strut\ttfamily\bfseries #1};%
  }%
}

\newcommand{\slicebody}[1]{%
  \begin{minipage}[t]{6.5cm}\ttfamily\tiny\raggedright #1\end{minipage}%
}
\newcommand{\rootbody}[1]{%
  \begin{minipage}[t]{13.5cm}\ttfamily\scriptsize\raggedright #1\end{minipage}%
}
\begin{tikzpicture}[
  codebox/.style={
    draw=black!55, line width=0.5pt, dotted,
    rounded corners=2.5pt, inner sep=5pt,
    anchor=north,
  },
  rootbox/.style={codebox, line width=0.7pt, draw=black!70, inner sep=6pt},
  artyarrow/.style={
    -{Stealth[length=2.8mm, width=2.2mm]},
    line width=0.7pt, color=black!55,
    decorate, decoration={
      snake, amplitude=0.6mm, segment length=2.7mm,
      pre length=2mm, post length=2.5mm,
    },
  },
  branchedge/.style={
    -{Stealth[length=2.5mm, width=2mm]},
    line width=0.6pt, color=black!55,
  },
  arrowlbl/.style={font=\itshape\scriptsize, color=black!60, fill=white, inner sep=1.5pt},
  stagelbl/.style={font=\itshape\footnotesize\color{black!55}, anchor=south, inner sep=1pt},
]

% --- ROOT: original program with red squiggle ---
% Standard-style syntax highlighting (not slice colours):
%   \kwd{...}   keywords (let, fun, type, case, ...) - muted purple
%   \ty{...}    types and type-level identifiers       - muted teal
%   default text for value-level identifiers and literals
\newcommand{\kwd}[1]{\textcolor{Purple!58!black}{#1}}
\newcommand{\ty}[1]{\textcolor{TealBlue!55!black}{#1}}
\node[rootbox] (orig) at (0,0) {\rootbody{%
\kwd{type} \ty{Option(A)} \kwd{=} \ty{Some(A)} \kwd{|} \ty{None}\\
\kwd{type} \ty{Digit} \kwd{=} \ty{Zero} \kwd{|} \ty{One} \kwd{|} \ty{Two} \kwd{|} $\cdots$ \kwd{|} \ty{Nine}\\
\kwd{type} \ty{Pin} \kwd{=} \ty{Digit × Digit × Digit × Digit}\\
\kwd{let} bind-parser \kwd{=} \kwd{typfun}\ty{<A>}\ \kwd{typfun}\ty{<B>}\\
\quad \kwd{fun}\ (p \kwd{:} \ty{String -\textgreater\ Option (String × A)}) \kwd{-\textgreater}\\
\quad \kwd{fun}\ (f \kwd{:} \ty{A -\textgreater\ String -\textgreater\ Option (String × B)}) \kwd{-\textgreater}\\
\quad \kwd{fun}\ (s \kwd{:} \ty{String}) \kwd{-\textgreater}\\
\quad\quad \kwd{case} p(s) \kwd{of}\\
\quad\quad\quad \kwd{|} \ty{None} \kwd{-\textgreater}\ \ty{None}\\
\quad\quad\quad \kwd{|} \ty{Some(}s', a\ty{)} \kwd{-\textgreater}\ f\ a\ s'\\
\kwd{let} digit\_parser \kwd{=} \kwd{fun}\ (s \kwd{:} \ty{String}) \kwd{-\textgreater}\\
\quad parse\_digit(s.head()).map(\kwd{fun}\ (d \kwd{:} \ty{Digit}) \kwd{-\textgreater}\ (s.tail(), d))\\
\kwd{let} parse\_pin \kwd{=} \kwd{fun}\ (s \kwd{:} \ty{String}) \kwd{-\textgreater}\\
\quad bind-parser\ty{<Digit><Pin>}\ digit\_parser (\kwd{fun}\ (d\textsubscript{1} \kwd{:} \ty{Digit}) \kwd{-\textgreater}\\
\quad bind-parser\ty{<Digit><Pin>}\ digit\_parser (\kwd{fun}\ (d\textsubscript{2} \kwd{:} \ty{Digit}) \kwd{-\textgreater}\\
\quad bind-parser\ty{<Digit><Pin>}\ digit\_parser (\kwd{fun}\ (d\textsubscript{3} \kwd{:} \ty{Digit}) \kwd{-\textgreater}\\
\quad\quad bind-parser\ty{<Digit><Pin>}\ digit\_parser (\\
\quad\quad\quad \focusbox{\squiggle{\kwd{fun}\ (d\textsubscript{4} \kwd{:} \ty{Digit}) \kwd{-\textgreater}\ \ty{Some(}s, (d\textsubscript{1}, d\textsubscript{2}, d\textsubscript{3}, d\textsubscript{4})\ty{)}}}))))) \\
\kwd{in}\ parse\_pin(\kwd{"1234"})%
}};

\node[stagelbl, anchor=south] at (orig.north) {original program};

% Error annotation: stretches to the full orig codebox width and sits flush
% against its bottom edge, with a dotted outline (matching the codebox style)
% and a very faint red wash.
\path let \p1=(orig.south west), \p2=(orig.south east) in
  node[anchor=north west, font=\scriptsize, align=center, inner sep=5pt,
       draw=red!70!black, line width=0.5pt, dotted, rounded corners=2.5pt,
       fill=red!4, text=red!55!black,
       text width={\x2-\x1-10pt}]
       (errbox) at (\p1)
       {term has type \synsrc{\texttt{Digit \textrightarrow\ Option (String × Pin)}}
        but was expected to have type \anasrc{\texttt{Digit \textrightarrow\ String \textrightarrow\ Option (String × Pin)}}};

% Column x-positions and explicit y-coordinates for each stage; using these
% rather than below-of chains guarantees both columns line up at every stage.
\def\colL{-4.05}
\def\colR{4.05}
% Y-coordinates of each box's NORTH anchor; queries sit ABOVE each box.
\coordinate (fullTop)    at (0,-11.8);
\coordinate (refinedTop) at (0,-19.3);

% Local: per-stage query label commands.
\newcommand{\synquery}[1]{%
  \begin{minipage}[t]{6.5cm}\centering
  {\scriptsize\itshape\color{black!55}query }{\scriptsize$\upsilon\,=\,$}{\scriptsize\synsrc{\texttt{#1}}}%
  \end{minipage}}
\newcommand{\anaquery}[1]{%
  \begin{minipage}[t]{6.5cm}\centering
  {\scriptsize\itshape\color{black!55}query }{\scriptsize$\upsilon\,=\,$}{\scriptsize\anasrc{\texttt{#1}}}%
  \end{minipage}}

% --- COLUMN: synthesis (left) ---
\node[codebox] (synFull) at ($(fullTop)+(\colL,0)$) {\slicebody{%
\strN{type} \asmref{Option}\synsrc{(A)} \strN{=} \synsrc{Some(A)} \strN{|} \fold\\
\strN{type} \asmref{Digit} \strN{=} \fold\\
\strN{type} \asmref{Pin} \strN{=} \synsrc{Digit × Digit × Digit × Digit}\\
\strN{let} \asmref{bind-parser} \strN{=} \fold\\
\fold\\
\strN{let} \fold\ \strN{=} \strN{fun} \synsrc{(}\asmref{s} \strN{:} \synsrcS{String}\synsrc{)} \strN{-\textgreater}\\
\quad \fold\ \fold\ \strN{(}\strN{fun} \synsrc{(}\asmref{d\textsubscript{1}} \strN{:} \synsrc{Digit}\synsrc{)} \strN{-\textgreater}\\
\quad \fold\ \fold\ \strN{(}\strN{fun} \synsrc{(}\asmref{d\textsubscript{2}} \strN{:} \synsrc{Digit}\synsrc{)} \strN{-\textgreater}\\
\quad \fold\ \fold\ \strN{(}\strN{fun} \synsrc{(}\asmref{d\textsubscript{3}} \strN{:} \synsrc{Digit}\synsrc{)} \strN{-\textgreater}\\
\quad \fold\ \fold\ \strN{(}\\
\quad\quad \focusbox{\squiggle{\synsrcS{fun}\ \synsrc{(}\asmref{d\textsubscript{4}} \synsrc{:} \synsrcS{Digit}\synsrc{)} \synsrcS{-\textgreater}\ \synsrcS{Some(}\asmref{s}\synsrcS{, (d\textsubscript{1}, d\textsubscript{2}, d\textsubscript{3}, d\textsubscript{4}))}}}\\
\quad \strN{))))}\\
\strN{in} \fold%
}};
\node[anchor=south, inner sep=3pt] (synFullQuery) at (synFull.north)
  {\synquery{Digit \textrightarrow\ Option (String × Pin)}};

\node[codebox] (synRefined) at ($(refinedTop)+(\colL,0)$) {\slicebody{%
\strN{type} \asmref{Option}\synsrc{(}\fold\synsrc{)} \strN{=} \synsrcS{Some(}\fold\synsrcS{)} \strN{|} \fold\\
\fold\\
\strN{let} \fold\ \strN{=} \strN{fun} \fold\ \strN{-\textgreater}\\
\quad \fold\ \fold\ \strN{(}\strN{fun} \fold\ \strN{-\textgreater}\\
\quad \fold\ \fold\ \strN{(}\strN{fun} \fold\ \strN{-\textgreater}\\
\quad \fold\ \fold\ \strN{(}\strN{fun} \fold\ \strN{-\textgreater}\\
\quad \fold\ \fold\ \strN{(}\\
\quad\quad \focusbox{\squiggle{\synsrcS{fun} \fold\ \synsrcS{-\textgreater}\ \synsrcS{Some} \fold}}\\
\quad \strN{))))}\\
\strN{in} \fold%
}};
\node[anchor=south, inner sep=3pt] (synRefinedQuery) at (synRefined.north)
  {\synquery{\fold\ \textrightarrow\ Option \fold}};

% --- COLUMN: analysis (right) ---
\node[codebox] (anaFull) at ($(fullTop)+(\colR,0)$) {\slicebody{%
\strN{type} \asmref{Option}\anasrc{(}\fold\anasrc{)} \strN{=} \fold\\
\strN{type} \asmref{Digit} \strN{=} \fold\\
\strN{type} \asmref{Pin} \strN{=} \fold\\
\strN{let} \asmref{bind-parser} \strN{=} \anasrc{typfun<A>} \anasrc{typfun<B>}\\
\quad \anasrc{fun} \anasrc{(}\fold\anasrc{)} \anasrc{-\textgreater}\\
\quad \anasrc{fun} \anasrc{(}\fold\ \strN{:} \anasrc{A} \anasrcS{-\textgreater\ String -\textgreater\ Option (String ×} \anasrc{B}\anasrcS{)}\anasrc{)} \anasrc{-\textgreater}\\
\quad \fold\\
\fold\\
\strN{let} \fold\ \strN{=} \strN{fun} \fold\ \strN{-\textgreater}\\
\quad \fold\ \fold\ \strN{(fun} \fold\ \strN{-\textgreater}\\
\quad \fold\ \fold\ \strN{(fun} \fold\ \strN{-\textgreater}\\
\quad \fold\ \fold\ \strN{(fun} \fold\ \strN{-\textgreater}\\
\quad \anasrc{bind-parser<}\anasrcS{Digit}\anasrc{><}\anasrcS{Pin}\anasrc{>}\ \fold\ \focusbox{\squiggle{\fold}}\strN{)))}\\
\strN{in} \fold%
}};
\node[anchor=south, inner sep=3pt] (anaFullQuery) at (anaFull.north)
  {\anaquery{Digit \textrightarrow\ String \textrightarrow\ Option (String × Pin)}};

\node[codebox] (anaRefined) at ($(refinedTop)+(\colR,0)$) {\slicebody{%
\fold\\
\strN{let} \asmref{bind-parser} \strN{=} \anasrc{typfun}\anasrc{<}\fold\anasrc{>} \anasrc{typfun}\anasrc{<}\fold\anasrc{>}\\
\quad \anasrc{fun} \anasrc{(}\fold\anasrc{)} \anasrc{-\textgreater}\\
\quad \anasrc{fun} \anasrc{(}\fold\ \strN{:} \fold\ \anasrcS{-\textgreater} \fold\ \anasrcS{-\textgreater} \fold\anasrc{)} \anasrc{-\textgreater}\\
\quad \fold\\
\strN{let} \fold\ \strN{=} \strN{fun} \fold\ \strN{-\textgreater}\\
\quad \fold\ \fold\ \strN{(fun} \fold\ \strN{-\textgreater}\\
\quad \fold\ \fold\ \strN{(fun} \fold\ \strN{-\textgreater}\\
\quad \fold\ \fold\ \strN{(fun} \fold\ \strN{-\textgreater}\\
\quad \anasrc{bind-parser}\anasrc{<}\fold\anasrc{><}\fold\anasrc{>}\ \fold\ \focusbox{\squiggle{\fold}}\strN{)))}\\
\strN{in} \fold%
}};
\node[anchor=south, inner sep=3pt] (anaRefinedQuery) at (anaRefined.north)
  {\anaquery{\fold\ \textrightarrow\ \fold\ \textrightarrow\ \fold}};

% --- Branching arrows from the root to each full-stage query ---
% Routed from the bottom of the error box so they curve around it.
\draw[branchedge] (errbox.south) to[bend right=22]
   node[arrowlbl, pos=0.5] {synthesis query} (synFullQuery.north);
\draw[branchedge] (errbox.south) to[bend left=22]
   node[arrowlbl, pos=0.5] {analysis query} (anaFullQuery.north);

% --- Inter-stage arrows: from each full box's south to the refined query above ---
\draw[artyarrow] (synFull.south) -- node[arrowlbl] {refine query} (synRefinedQuery.north);
\draw[artyarrow] (anaFull.south) -- node[arrowlbl] {refine query} (anaRefinedQuery.north);

% Stage labels in the gutter between the two columns
\node[stagelbl, color=black!55, anchor=center] at (synFull.center    -| 0,0) {full};
\node[stagelbl, color=black!55, anchor=center] at (synRefined.center -| 0,0) {refined};

\end{tikzpicture}

\caption{Debugging a static type error in a parser for 4 digit PINs.}
\label{fig:parse-pin-slices}
\end{figure}

\begin{figure}[p]
\centering

% Local visual conventions (figure-scoped). Reused from previous figures.
\providecommand{\fold}{%
  \tikz[baseline=0.1em]{%
    \draw[rounded corners=0.3pt, line width=0.28pt, black!75]
      (0,0) rectangle (0.85em,0.85em);%
    \fill[black!85] (0.18em,0.67em) circle (0.07em);%
    \fill[black!85] (0.425em,0.425em) circle (0.07em);%
    \fill[black!85] (0.67em,0.18em) circle (0.07em);%
  }%
}
\providecommand{\synsrc}[1]{\textcolor{BrickRed!60!gray}{#1}}
\providecommand{\anasrc}[1]{\textcolor{NavyBlue!60!gray}{#1}}
\providecommand{\synsrcS}[1]{\textcolor{red!85!black}{\textbf{\boldmath\addfontfeature{FakeBold=4.5}#1}}}
\providecommand{\anasrcS}[1]{\textcolor{blue!85!black}{\textbf{\boldmath\addfontfeature{FakeBold=4.5}#1}}}
\providecommand{\strN}[1]{\textcolor{OliveGreen!65!black}{#1}}
\providecommand{\strM}[1]{\textcolor{OliveGreen!65!black}{#1}}
\providecommand{\strF}[1]{\textcolor{OliveGreen!65!black}{#1}}
\providecommand{\offslice}[1]{\textcolor{black!42}{#1}}
\providecommand{\focusbox}[1]{%
  \tikz[baseline=(N.base)]{%
    \node[draw=black!30, line width=0.3pt, rounded corners=2pt,
          inner sep=2pt, fill=black!10, anchor=base, align=left] (N)
         {\strut\ttfamily\bfseries #1};%
  }%
}

\newcommand{\slicebody}[1]{%
  \begin{minipage}[t]{6.5cm}\ttfamily\tiny\raggedright #1\end{minipage}%
}
\newcommand{\fullslicebody}[1]{%
  \begin{minipage}[t]{12cm}\ttfamily\tiny\raggedright #1\end{minipage}%
}
\newcommand{\rootbody}[1]{%
  \begin{minipage}[t]{13.5cm}\ttfamily\scriptsize\raggedright #1\end{minipage}%
}

\begin{tikzpicture}[
  codebox/.style={
    draw=black!55, line width=0.5pt, dotted,
    rounded corners=2.5pt, inner sep=5pt,
    anchor=north,
  },
  rootbox/.style={codebox, line width=0.7pt, draw=black!70, inner sep=6pt},
  artyarrow/.style={
    -{Stealth[length=2.8mm, width=2.2mm]},
    line width=0.7pt, color=black!55,
    decorate, decoration={
      snake, amplitude=0.6mm, segment length=2.7mm,
      pre length=2mm, post length=2.5mm,
    },
  },
  branchedge/.style={
    -{Stealth[length=2.5mm, width=2mm]},
    line width=0.6pt, color=black!55,
  },
  arrowlbl/.style={font=\itshape\scriptsize, color=black!60, fill=white, inner sep=1.5pt},
  stagelbl/.style={font=\itshape\footnotesize\color{black!55}, anchor=south, inner sep=1pt},
]

% --- ROOT: original program. Syntax highlighting (NOT slice colours): purple
%   keywords (\kwd), teal types (\ty), default identifiers and literals.
\newcommand{\kwd}[1]{\textcolor{Purple!58!black}{#1}}
\newcommand{\ty}[1]{\textcolor{TealBlue!55!black}{#1}}
\node[rootbox] (orig) at (0,0) {\rootbody{%
\kwd{type} \ty{Option(A)} \kwd{=} \ty{Some(A)} \kwd{|} \ty{None}\\
\kwd{type} \ty{Digit} \kwd{=} \ty{Zero} \kwd{|} \ty{One} \kwd{|} \ty{Two} \kwd{|} $\cdots$ \kwd{|} \ty{Nine}\\
\kwd{let} parse\_digit \kwd{=} \kwd{fun}\ (c \kwd{:} \ty{Char}) \kwd{-\textgreater}\ \kwd{case}\ c\ \kwd{of}\\
\quad \kwd{|} \kwd{'0'} \kwd{=\textgreater}\ \ty{Some Zero} \kwd{|} \kwd{'1'} \kwd{=\textgreater}\ \ty{Some One} \kwd{|} $\cdots$ \kwd{|} \kwd{'9'} \kwd{=\textgreater}\ \ty{Some Nine} \kwd{|} \_ \kwd{=\textgreater}\ \ty{None}\\
\kwd{let} parse\_string \kwd{:} \ty{Char -\textgreater\ String} \kwd{=} string\_of\_char\\
\kwd{let} lex\_atom \kwd{=} \kwd{fun}\ (c \kwd{:} \ty{Char}) \kwd{-\textgreater}\\
\quad \kwd{if} is\_digit(c) \kwd{then} \ty{Left(}parse\_digit(c).get()\ty{)} \kwd{else} \ty{Right(}parse\_string(c)\ty{)}\\
\kwd{in}\ \focusbox{lex\_atom}(\kwd{`5'})%
}};

\node[stagelbl, anchor=south] at (orig.north) {original program};

% Column x-positions and explicit y-coordinates for each stage.
\def\colC{0}
\def\colL{-4.05}
\def\colR{4.05}
\coordinate (fullTop)    at (\colC,-4.7);
\coordinate (refinedTop) at (0,-10.6);

% Local: per-stage query label commands. (Reused style from figure 2.)
\newcommand{\synqueryT}[1]{%
  \begin{minipage}[t]{6.5cm}\centering
  {\scriptsize\itshape\color{black!55}query }{\scriptsize$\upsilon\,=\,$}{\scriptsize\synsrc{\texttt{#1}}}%
  \end{minipage}}

% --- FULL synthesis slice (centred under root) ---
\node[codebox] (fullSlice) at (fullTop) {\fullslicebody{%
\strN{type} \asmref{Option}\synsrc{(A)} \strN{=} \synsrc{Some(A)} \strN{|} \fold\\
\strN{type} \asmref{Digit} \strN{=} \synsrcS{Zero} \strN{|} \fold\\
\strN{let} \asmref{parse\_digit} \strN{=} \synsrc{fun} \fold\ \synsrc{-\textgreater}\ \strN{case} \fold\ \strN{of}\\
\quad \strF{|} \fold\ \strF{=\textgreater} \synsrc{Some} \synsrcS{Zero} \strF{|} \fold\\
\strN{let} \asmref{parse\_string} \strN{:} \fold\ \synsrc{-\textgreater\ String} \strN{=} \fold\\
\strN{let} \asmref{lex\_atom} \strN{=} \synsrcS{fun} \synsrc{(}\fold\ \strN{:} \synsrcS{Char}\synsrc{)} \synsrcS{-\textgreater}\\
\quad \synsrcS{if} \fold\ \synsrc{then} \synsrcS{Left}\synsrc{(}\synsrc{parse\_digit}\synsrc{(}\fold\synsrc{)}\synsrc{.get()}\synsrc{)} \synsrc{else} \synsrcS{Right}\synsrc{(}\synsrc{parse\_string}\synsrc{(}\fold\synsrc{)}\synsrc{)}\\
\strN{in} \focusbox{\synsrc{lex\_atom}}\strN{(}\fold\strN{)}%
}};
\node[anchor=south, inner sep=3pt] (fullQuery) at (fullSlice.north)
  {\synqueryT{Char \textrightarrow\ Digit + String}};

% --- LEFT refined slice: option-digit half ---
\node[codebox] (legSlice) at ($(refinedTop)+(\colL,0)$) {\slicebody{%
\strN{type} \asmref{Option}\synsrc{(A)} \strN{=} \synsrc{Some(A)} \strN{|} \fold\\
\strN{type} \asmref{Digit} \strN{=} \synsrcS{Zero} \strN{|} \fold\\
\strN{let} \asmref{parse\_digit} \strN{=} \synsrc{fun} \fold\ \synsrc{-\textgreater}\ \strN{case} \fold\ \strN{of}\\
\quad \strF{|} \fold\ \strF{=\textgreater} \synsrc{Some} \synsrcS{Zero} \strF{|} \fold\\
\strN{let} \asmref{lex\_atom} \strN{=} \synsrcS{fun} \synsrc{(}\fold\ \strN{:} \synsrcS{Char}\synsrc{)} \synsrcS{-\textgreater}\\
\quad \strN{if} \fold\ \strN{then} \synsrcS{Left}\synsrc{(}\synsrc{parse\_digit}\synsrc{(}\fold\synsrc{)}\synsrc{.get()}\synsrc{)} \strN{else} \fold\\
\strN{in} \focusbox{\synsrc{lex\_atom}}\strN{(}\fold\strN{)}%
}};
\node[anchor=south, inner sep=3pt] (legQuery) at (legSlice.north)
  {\synqueryT{Char \textrightarrow\ Digit + \fold}};

% --- RIGHT refined slice: string half ---
\node[codebox] (restSlice) at ($(refinedTop)+(\colR,0)$) {\slicebody{%
\strN{let} \asmref{parse\_string} \strN{:} \fold\ \synsrc{-\textgreater\ String} \strN{=} \fold\\
\strN{let} \asmref{lex\_atom} \strN{=} \synsrcS{fun} \synsrc{(}\fold\ \strN{:} \synsrcS{Char}\synsrc{)} \synsrcS{-\textgreater}\\
\quad \strN{if} \fold\ \strN{then} \fold\ \strN{else} \synsrcS{Right}\synsrc{(}\synsrc{parse\_string}\synsrc{(}\fold\synsrc{)}\synsrc{)}\\
\strN{in} \focusbox{\synsrc{lex\_atom}}\strN{(}\fold\strN{)}%
}};
\node[anchor=south, inner sep=3pt] (restQuery) at (restSlice.north)
  {\synqueryT{Char \textrightarrow\ \fold\ + String}};

% --- Arrows ---
% Root -> full slice (centred)
\draw[branchedge] (orig.south) -- node[arrowlbl] {synthesis query} (fullQuery.north);

% Full slice -> two refined slices (fan out)
\draw[artyarrow] (fullSlice.south) to[bend right=18]
   node[arrowlbl, pos=0.55] {refine to digit half} (legQuery.north);
\draw[artyarrow] (fullSlice.south) to[bend left=18]
   node[arrowlbl, pos=0.55] {refine to string half} (restQuery.north);

\end{tikzpicture}

\caption{Lexing a character into a \texttt{Left(Digit)} or \texttt{Right(String)} token.}
\label{fig:parse-event-slices}
\end{figure}

%%% Local Variables:
%%% mode: LaTeX
%%% TeX-master: "PolymorphicTypeSlicing"
%%% End:
